% Axial resolution from parallax, drawn in object space without lenses
% (supervisor, 2026-11-17 meeting, workplan D9). Two on-axis emitters, the
% parallel lines of sight of the outermost view at theta_max, the focal plane,
% and the displacement Delta z tan(theta_max) compared with Delta x_LF.
% Since 2026-11-23 (supervisor meeting) theta_max is the pupil edge,
% sin(theta_max) ~ NA = 0.75 (was the center of the outermost view, 0.6);
% the ray length and Delta z were shortened to keep the height.
% Not to scale. The lens version is kept in axial-parallax-lenses.tex
% (Appendix D).
\begin{tikzpicture}[
    x=1cm,y=1cm,
    font=\footnotesize,
    ink/.style={text=black!85},
    rayA/.style={draw=blue!60!black,line width=0.8pt},
    rayB/.style={draw=red!65!black,line width=0.8pt,dash pattern=on 3pt off 2pt},
    plane/.style={draw=black!45,dashed,line width=0.5pt},
    dim/.style={<->,>={Stealth[length=1.6mm]},draw=black!70,line width=0.5pt},
    lead/.style={draw=black!55,line width=0.4pt},
    lab/.style={font=\footnotesize,ink,inner sep=1.5pt,align=center}
]
% ---- geometry (cm; not to scale) ----
\def\dz{1.6}                              % axial separation, exaggerated
\def\th{48.59}                            % theta_max, sin = 0.75 = NA (pupil edge)
\pgfmathsetmacro{\tn}{tan(\th)}
\pgfmathsetmacro{\sh}{\dz*\tn}            % displacement in the focal plane
\def\xe{2.6}                              % where the drawn rays end
\pgfmathsetmacro{\hA}{\xe*\tn}
\pgfmathsetmacro{\hB}{(\xe+\dz)*\tn}

% optical axis = line of sight of the central view
\draw[->,>={Stealth[length=2mm]},draw=green!45!black,line width=0.9pt]
    (-\dz-0.8,0) -- (\xe+2.2,0);
\node[lab,anchor=north,text=green!45!black] at (\xe+0.5,-0.15)
    {central view:\\both emitters coincide};

% focal plane
\draw[plane] (0,-1.3) -- (0,\hB+0.2);
\node[lab,anchor=south] at (0,\hB+0.2) {focal plane};

% lines of sight of the outermost view (parallel, at theta_max)
\draw[rayA,->,>={Stealth[length=2mm]}] (0,0) -- (\xe,\hA);
\draw[rayB,->,>={Stealth[length=2mm]}] (-\dz,0) -- (\xe,\hB);
\node[lab,anchor=west] at (\xe+0.1,\hA+0.75) {to the\\outermost view};

% angle theta_max
\draw[draw=black!70,line width=0.5pt] (1.3,0) arc (0:\th:1.3);
\node[lab,anchor=west] at (1.35,0.42) {$\theta_{max}$};

% emitters and the apparent position of the defocused one
\fill[blue!60!black] (0,0) circle (0.1);
\fill[red!65!black] (-\dz,0) circle (0.1);
\draw[red!65!black,line width=0.7pt,fill=white] (0,\sh) circle (0.09);
\node[lab,anchor=north west,text=blue!60!black,fill=white] at (0.08,-0.12) {in focus};
\node[lab,anchor=north,text=red!65!black,fill=white] at (-\dz,-0.14) {out of focus};

% displacement in the focal plane
\draw[dim] (0.3,0.02) -- (0.3,\sh-0.02);
% label in the clear wedge right of the in-focus ray, above the theta arc
\draw[lead] (0.34,\sh/2+0.3) -- (1.75,\sh/2+0.6);
\node[lab,anchor=west] at (1.75,\sh/2+0.6) {$\Delta z\tan\theta_{max}$};

% axial separation
\draw[dim] (-\dz,-0.85) -- (0,-0.85);
\node[lab,anchor=north] at (-\dz/2,-0.9) {$\Delta z$};

% ---- right panel: the two images in the outermost view at the limit ----
\def\xi{7.6}
\def\ya{0.6}                              % image of the focused emitter
\def\yb{2.0}                              % image of the defocused emitter
\draw[->,>={Stealth[length=1.6mm]},draw=black!70,line width=0.5pt]
    (\xi-0.3,-0.9) -- (\xi-0.3,3.7);
\node[lab,anchor=east] at (\xi-0.45,3.3) {lateral};
\draw[draw=blue!60!black,line width=0.8pt,samples=60,domain=\ya-1.5:\ya+1.5]
    plot ({\xi+1.6*exp(-2.7726*((\x-\ya)/1.4)^2)},{\x});
\draw[draw=red!65!black,line width=0.8pt,dash pattern=on 3pt off 2pt,
    samples=60,domain=\yb-1.5:\yb+1.5]
    plot ({\xi+1.6*exp(-2.7726*((\x-\yb)/1.4)^2)},{\x});
\draw[dim] (\xi+1.95,\ya) -- (\xi+1.95,\yb);
\node[lab,anchor=west] at (\xi+2.0,{(\ya+\yb)/2}) {$\Delta x_{LF}$};
\node[lab,anchor=south] at (\xi+0.9,3.8) {outermost view,\\at the limit};
\end{tikzpicture}
